% Lines 1143-1597 from original attention_transformers_notes_ghost.tex
% Chapter 6: The Transformer Architecture

%==============================================================================
\section{The Transformer Architecture}
%==============================================================================

\subsection{Overview: Attention is All You Need}

The transformer architecture, introduced by Vaswani et al. (2017), revolutionized sequence modeling by relying entirely on attention mechanisms, dispensing with recurrence and convolutions. The key insights were:

\begin{enumerate}
    \item \textbf{Parallel Processing}: All positions can be processed simultaneously
    \item \textbf{Direct Connections}: Any two positions can interact directly through attention
    \item \textbf{Computational Efficiency}: Despite $O(n^2)$ attention complexity, parallelization makes training much faster
    \item \textbf{Interpretability}: Attention weights provide insight into model decisions
\end{enumerate}

\subsection{Architecture Components}

The transformer consists of an encoder and decoder, each composed of stacked identical layers.

\begin{figure}[H]
\centering
\tikzset{
    enc/.style={draw, fill=nodegreen!45, minimum width=3.4cm, minimum height=0.9cm, rounded corners=2pt},
    dec/.style={draw, fill=nodeblue!45, minimum width=3.4cm, minimum height=0.9cm, rounded corners=2pt},
    norm/.style={draw, fill=gray!30, minimum width=3.4cm, minimum height=0.55cm, rounded corners=2pt},
    ff/.style={draw, fill=nodegreen!20, minimum width=3.4cm, minimum height=0.9cm, rounded corners=2pt},
    box/.style={draw, fill=gray!15, minimum width=3.2cm, minimum height=0.8cm, rounded corners=2pt},
    thinbox/.style={draw, fill=yellow!20, minimum width=3.2cm, minimum height=0.8cm, rounded corners=2pt},
    label/.style={font=\bfseries},
    arr/.style={->,>=stealth,thick},
    res/.style={->,>=stealth,thick,rounded corners=2pt,black!70}
}
\begin{tikzpicture}[node distance=7mm and 18mm]
    % Bottom: embeddings + positional encoding + add
    % Left (Encoder side)
    \node[box] (embE) at (-6,-7) {Input Embedding};
    \node[thinbox, above=4mm of embE] (posE) {Positional Encoding};
    \node[circle, draw, above=4mm of posE, minimum size=6mm] (addE) {+};
    \draw[arr] (embE) -- (posE);
    \draw[arr] (posE) -- (addE);

    % Right (Decoder side)
    \node[box] (embD) at (6,-7) {Output Embedding (shifted)};
    \node[thinbox, above=4mm of embD] (posD) {Positional Encoding};
    \node[circle, draw, above=4mm of posD, minimum size=6mm] (addD) {+};
    \draw[arr] (embD) -- (posD);
    \draw[arr] (posD) -- (addD);

    % Encoder stack (single layer, bracketed N×)
    \node[enc, above=14mm of addE] (e_att1) {Multi-Head Attention};
    \node[norm, above=3mm of e_att1] (e_norm1) {Add \& Norm};
    \node[ff, above=3mm of e_norm1] (e_ff1) {Feed Forward};
    \node[norm, above=3mm of e_ff1] (e_norm2) {Add \& Norm};
    % Entry to first encoder sublayer via internal anchor (keeps skip inside box)
    \path (e_att1.west) ++(-0.25,0) coordinate (e_in);
    \draw[arr] (addE) -- (e_in) -- (e_att1.west);
    \draw[arr] (e_att1) -- (e_norm1);
    \draw[arr] (e_norm1) -- (e_ff1);
    \draw[arr] (e_ff1) -- (e_norm2);

    % Encoder residual (skip) connections — drawn inside bracket
    \draw[res] (e_in) -- ++(-0.5,0) |- (e_norm1.west);
    \draw[res] (e_norm1.west) -- ++(-0.5,0) |- (e_norm2.west);

    % Encoder bracket N×
    \node[draw, dashed, inner sep=12pt, fit=(e_att1)(e_norm2)] (encbox) {};
    \node[below left=0.1cm of encbox] {N×};
    \node[above=2mm of encbox, label] {Encoder};

    % Decoder stack (one layer, bracketed N×): masked self-attn -> norm -> cross-attn -> norm -> FF -> norm
    \node[dec, above=14mm of addD] (d_self) {Masked Multi-Head Attention};
    \node[norm, above=3mm of d_self] (d_norm1) {Add \& Norm};
    \node[dec, above=3mm of d_norm1] (d_cross) {Multi-Head Attention (cross)};
    \node[norm, above=3mm of d_cross] (d_norm2) {Add \& Norm};
    \node[dec, above=3mm of d_norm2, fill=nodeblue!35] (d_ff) {Feed Forward};
    \node[norm, above=3mm of d_ff] (d_norm3) {Add \& Norm};
    % Entry to first decoder sublayer via internal anchor
    \path (d_self.west) ++(-0.25,0) coordinate (d_in);
    \draw[arr] (addD) -- (d_in) -- (d_self.west);
    \draw[arr] (d_self) -- (d_norm1);
    \draw[arr] (d_norm1) -- (d_cross);
    \draw[arr] (d_cross) -- (d_norm2);
    \draw[arr] (d_norm2) -- (d_ff);
    \draw[arr] (d_ff) -- (d_norm3);

    % Decoder residual (skip) connections — inside bracket
    \draw[res] (d_in) -- ++(-0.5,0) |- (d_norm1.west);
    \draw[res] (d_norm1.west) -- ++(-0.5,0) |- (d_norm2.west);
    \draw[res] (d_norm2.west) -- ++(-0.5,0) |- (d_norm3.west);

    % Decoder head
    \node[box, above=12mm of d_norm3] (linear) {Linear};
    \node[box, above=3mm of linear, fill=orange!20] (softmax) {Softmax};
    \node[box, above=3mm of softmax] (outprob) {Output Probabilities};
    \draw[arr] (d_norm3) -- (linear);
    \draw[arr] (linear) -- (softmax);
    \draw[arr] (softmax) -- (outprob);

    % Decoder bracket N×
    \node[draw, dashed, inner sep=12pt, fit=(d_self)(d_norm3)] (decbox) {};
    \node[below left=0.1cm of decbox] {N×};
    \node[above=2mm of decbox, label] {Decoder};

    % Cross-attention connections (clean, curved)
    \draw[arr, red] (e_norm2.east) to[bend left=12] node[midway, above, sloped] {K,V} (d_cross.west);

    % Titles
    \node[font=\Large\bfseries, above=10mm of outprob] {The Transformer Architecture};
\end{tikzpicture}
\caption{Complete Transformer Architecture}
\label{fig:transformer-architecture}
\end{figure}

\subsection{Encoder Layer}

Each encoder layer consists of two sub-layers:
\begin{enumerate}
    \item Multi-head self-attention
    \item Position-wise feed-forward network
\end{enumerate}

\subsubsection{Residual Connections and Layer Normalization}

Around each sub-layer, we employ:
\begin{itemize}
    \item \textbf{Residual connection}: Helps with gradient flow and enables deep networks
    \item \textbf{Layer normalization}: Stabilizes training
\end{itemize}

The output of each sub-layer is:
\begin{equation}
\text{LayerNorm}(x + \text{Sublayer}(x))
\end{equation}
\begin{notationbox}
\textbf{Notation:}
\begin{itemize}
    \item $x$: Input to the sub-layer
    \item $\text{Sublayer}(x)$: The sub-layer function (attention or feed-forward)
    \item $\text{LayerNorm}$: Layer normalization
\end{itemize}
\end{notationbox}

\subsubsection{Position-wise Feed-Forward Networks}

The feed-forward network applies the same fully connected network to each position:

\begin{equation}
\text{FFN}(x) = \max(0, x\mathbf{W}_1 + \mathbf{b}_1)\mathbf{W}_2 + \mathbf{b}_2
\end{equation}
\begin{notationbox}
\textbf{Notation:}
\begin{itemize}
    \item $\mathbf{W}_1 \in \mathbb{R}^{d_{model} \times d_{ff}}$: First layer weights
    \item $\mathbf{W}_2 \in \mathbb{R}^{d_{ff} \times d_{model}}$: Second layer weights
    \item $d_{ff}$: Feed-forward dimension (typically $4 \times d_{model}$)
    \item ReLU activation between layers
\end{itemize}
\end{notationbox}

\subsection{Decoder Layer}

Each decoder layer has three sub-layers:
\begin{enumerate}
    \item Masked multi-head self-attention
    \item Multi-head cross-attention (encoder-decoder attention)
    \item Position-wise feed-forward network
\end{enumerate}

\subsubsection{Masked Self-Attention}

The decoder uses masked self-attention to prevent positions from attending to subsequent positions. This ensures the model can only use information from past positions during training, maintaining the autoregressive property.

The mask is implemented by setting attention scores to $-\infty$ before softmax:
\begin{equation}
\text{Mask}_{ij} = \begin{cases}
0 & \text{if } i \geq j \\
-\infty & \text{if } i < j
\end{cases}
\end{equation}

\subsubsection{Cross-Attention}

The second attention layer in the decoder performs multi-head attention over the encoder output:
\begin{itemize}
    \item Queries come from the previous decoder layer
    \item Keys and values come from the encoder output
    \item This allows every decoder position to attend to all encoder positions
\end{itemize}

\subsection{Positional Encoding}

Since transformers have no inherent notion of position or order, we must inject positional information. The original paper uses sinusoidal positional encodings:

\begin{align}
PE_{(pos,2i)} &= \sin(pos/10000^{2i/d_{model}}) \\
PE_{(pos,2i+1)} &= \cos(pos/10000^{2i/d_{model}})
\end{align}
\begin{notationbox}
\textbf{Notation:}
\begin{itemize}
    \item $pos$: Position in the sequence
    \item $i$: Dimension index
    \item $d_{model}$: Model dimension
    \item Each dimension uses a different frequency
\end{itemize}
\end{notationbox}

\subsubsection{Why Sinusoidal?}

Sinusoidal encodings have several advantages:
\begin{enumerate}
    \item \textbf{Unique encoding}: Each position gets a unique encoding
    \item \textbf{Relative positions}: The encoding for position $pos + k$ can be represented as a linear function of the encoding at $pos$
    \item \textbf{Extrapolation}: Can handle sequences longer than those seen during training
    \item \textbf{No parameters}: Unlike learned positional embeddings
\end{enumerate}

\subsection{Theoretical Analysis of Positional Encodings}

\subsubsection{Fourier Analysis of Sinusoidal Encodings}

The sinusoidal positional encoding can be understood through Fourier analysis, revealing deep connections to signal processing.

\textbf{Theorem (Fourier Basis Interpretation):} The sinusoidal positional encoding forms a truncated Fourier basis where each position is represented as:
\begin{equation}
\text{PE}(pos) = \sum_{i=0}^{d_{model}/2-1} \left[\sin(\omega_i \cdot pos) \cdot \mathbf{e}_{2i} + \cos(\omega_i \cdot pos) \cdot \mathbf{e}_{2i+1}\right]
\end{equation}
where $\omega_i = 1/10000^{2i/d_{model}}$ are the frequencies.

\begin{notationbox}
\textbf{Notation:}
\begin{itemize}
    \item $\omega_i$: Frequency for dimension $i$
    \item $\mathbf{e}_j$: Standard basis vector for dimension $j$
    \item The frequencies form a geometric progression
\end{itemize}
\end{notationbox}

\textbf{Proposition (Relative Position Property):} For any fixed offset $k$, there exists a linear transformation $\mathbf{T}_k$ such that:
\begin{equation}
\text{PE}(pos + k) = \mathbf{T}_k \cdot \text{PE}(pos)
\end{equation}

\textbf{Proof:} Using trigonometric identities:
\begin{align}
\sin(\omega(pos + k)) &= \sin(\omega \cdot pos)\cos(\omega \cdot k) + \cos(\omega \cdot pos)\sin(\omega \cdot k) \\
\cos(\omega(pos + k)) &= \cos(\omega \cdot pos)\cos(\omega \cdot k) - \sin(\omega \cdot pos)\sin(\omega \cdot k)
\end{align}

This can be written in matrix form where $\mathbf{T}_k$ is a block-diagonal matrix with 2×2 rotation matrices.

\subsubsection{Information-Theoretic Analysis of Position Representations}

We analyze how much positional information different encoding schemes can represent.

\textbf{Definition (Positional Information Capacity):} For encoding function $\text{PE}: \mathbb{N} \rightarrow \mathbb{R}^d$:
\begin{equation}
\mathcal{C}(\text{PE}) = \lim_{n \rightarrow \infty} \frac{\log |\{\text{PE}(i) : i \in [n]\}|}{n}
\end{equation}

\textbf{Caveat (Sinusoidal capacity):} Sinusoidal encodings do not add parameters with length, and pairwise distances can become arbitrarily small as positions grow. They still support linear relative-position computation, but there is no positive lower bound on $\|\text{PE}(i)-\text{PE}(j)\|_2$ as $i,j \to \infty$.

\subsubsection{Learned vs Fixed Positional Encodings: Theoretical Trade-offs}

\textbf{Theorem (Approximation Power):} Any continuous position-dependent function $f: [n] \times \mathbb{R}^d \rightarrow \mathbb{R}^d$ can be approximated by:
\begin{itemize}
    \item Learned embeddings: exactly (with $n$ parameters)
    \item Sinusoidal: within error $\epsilon$ (with $O(\log(n/\epsilon))$ dimensions)
\end{itemize}

\textbf{Proof Sketch:} Learned embeddings can memorize any function on finite positions. Sinusoidal encodings approximate via Fourier series with convergence rate dependent on function smoothness.

\subsubsection{Alternative Positional Encoding Schemes}

We analyze theoretical properties of various positional encoding schemes:

\textbf{1. Relative Positional Encoding:}
\begin{equation}
\text{Attention}_{ij} = \frac{\exp((\mathbf{q}_i^T\mathbf{k}_j + a_{i-j})/\sqrt{d_k})}{\sum_k \exp((\mathbf{q}_i^T\mathbf{k}_k + a_{i-k})/\sqrt{d_k})}
\end{equation}

where $a_{i-j}$ is a learned bias based on relative position.

\textbf{Property:} Translation invariant - shifting all positions preserves relative attention patterns.

\textbf{2. Rotary Position Embedding (RoPE):}
\begin{equation}
f_{\text{RoPE}}(\mathbf{x}, m) = \mathbf{R}_{\Theta,m} \mathbf{x}
\end{equation}

where $\mathbf{R}_{\Theta,m}$ is a rotation matrix dependent on position $m$.

\textbf{Property:} Preserves inner products up to position-dependent rotation:
\begin{equation}
\langle f(\mathbf{q}, m), f(\mathbf{k}, n) \rangle = \langle \mathbf{q}, \mathbf{R}_{\Theta,n-m} \mathbf{k} \rangle
\end{equation}

\subsubsection{Theoretical Limitations of Position-Agnostic Attention}

\textbf{Theorem (Permutation Invariance):} Without positional encoding, self-attention is permutation equivariant:
\begin{equation}
\text{Attention}(\mathbf{P}\mathbf{X}) = \mathbf{P}\text{Attention}(\mathbf{X})
\end{equation}
for any permutation matrix $\mathbf{P}$.

\textbf{Corollary:} Position-agnostic transformers cannot solve tasks requiring position-dependent computation (e.g., sorting, counting).

\textbf{Proof:} Any position-dependent function $f$ satisfies $f(\mathbf{X}) \neq f(\mathbf{P}\mathbf{X})$ for some permutation $\mathbf{P}$, but attention without positional encoding cannot distinguish these inputs.

\subsection{Complete Transformer Implementation}

Let's implement the complete transformer encoder:

\begin{lstlisting}[language=Python, caption=Transformer Encoder Implementation]
from __future__ import annotations
from dataclasses import dataclass
import torch
import torch.nn as nn
import torch.nn.functional as F
from typing import Optional
import math


@dataclass(frozen=True, slots=True)
class TransformerConfig:
    d_model: int = 512
    n_heads: int = 8
    n_layers: int = 6
    d_ff: int = 2048
    dropout: float = 0.1
    max_seq_len: int = 5000

    def __post_init__(self) -> None:
        if self.d_model % self.n_heads != 0:
            raise ValueError("d_model must be divisible by n_heads")
        if self.d_ff < self.d_model:
            raise ValueError("d_ff should typically be larger than d_model")


class PositionalEncoding(nn.Module):
    def __init__(self, config: TransformerConfig) -> None:
        super().__init__()
        self._d_model: int = config.d_model
        self._dropout: nn.Dropout = nn.Dropout(config.dropout)

        pe: torch.Tensor = self._create_positional_encoding(
            config.max_seq_len, config.d_model
        )
        self.register_buffer('pe', pe)

    def _create_positional_encoding(
        self,
        max_len: int,
        d_model: int
    ) -> torch.Tensor:
        pe: torch.Tensor = torch.zeros(max_len, d_model)
        position: torch.Tensor = torch.arange(0, max_len).unsqueeze(1).float()

        div_term: torch.Tensor = torch.exp(
            torch.arange(0, d_model, 2).float() *
            -(math.log(10000.0) / d_model)
        )

        pe[:, 0::2] = torch.sin(position * div_term)
        pe[:, 1::2] = torch.cos(position * div_term)

        return pe.unsqueeze(0)

    def forward(self, x: torch.Tensor) -> torch.Tensor:
        seq_len: int = x.size(1)
        x = x + self.pe[:, :seq_len, :]
        return self._dropout(x)


class TransformerEncoderLayer(nn.Module):
    def __init__(self, config: TransformerConfig) -> None:
        super().__init__()
        self._config: TransformerConfig = config

        attention_config: AttentionConfig = AttentionConfig(
            d_model=config.d_model,
            n_heads=config.n_heads,
            dropout=config.dropout
        )

        self._self_attention: MultiHeadAttention = MultiHeadAttention(
            attention_config
        )
        self._feed_forward: FeedForwardNetwork = FeedForwardNetwork(
            config.d_model, config.d_ff, config.dropout
        )

        self._norm1: nn.LayerNorm = nn.LayerNorm(config.d_model)
        self._norm2: nn.LayerNorm = nn.LayerNorm(config.d_model)

        self._dropout: nn.Dropout = nn.Dropout(config.dropout)

    def forward(
        self,
        x: torch.Tensor,
        mask: Optional[torch.Tensor] = None
    ) -> torch.Tensor:
        # Self-attention with residual connection
        attn_output: torch.Tensor = self._self_attention(x, x, x, mask)
        x = self._norm1(x + self._dropout(attn_output))

        # Feed-forward with residual connection
        ff_output: torch.Tensor = self._feed_forward(x)
        x = self._norm2(x + self._dropout(ff_output))

        return x


class FeedForwardNetwork(nn.Module):
    def __init__(
        self,
        d_model: int,
        d_ff: int,
        dropout: float = 0.1
    ) -> None:
        super().__init__()
        self._linear1: nn.Linear = nn.Linear(d_model, d_ff)
        self._linear2: nn.Linear = nn.Linear(d_ff, d_model)
        self._dropout: nn.Dropout = nn.Dropout(dropout)

    def forward(self, x: torch.Tensor) -> torch.Tensor:
        x = F.relu(self._linear1(x))
        x = self._dropout(x)
        x = self._linear2(x)
        return x


class TransformerEncoder(nn.Module):
    def __init__(self, config: TransformerConfig) -> None:
        super().__init__()
        self._config: TransformerConfig = config

        self._layers: nn.ModuleList = nn.ModuleList([
            TransformerEncoderLayer(config)
            for _ in range(config.n_layers)
        ])

        self._norm: nn.LayerNorm = nn.LayerNorm(config.d_model)

    def forward(
        self,
        x: torch.Tensor,
        mask: Optional[torch.Tensor] = None
    ) -> torch.Tensor:
        for layer in self._layers:
            x = layer(x, mask)

        return self._norm(x)
\end{lstlisting}
